\section{实验}

\subsection{环境设置}

实验采用二维异构无人机协同追逃环境，包含 3 架追击无人机和 1 个机动目标。追击无人机具有不同最大速度、感知范围和能耗系数。目标采用 mixed 机动策略，包括远离最近追击者逃逸和随机转向机动。最终主结果采用 mixed 目标、目标速度 0.75、通信半径 $4,6,8,10$、每种子 300 回合评估和 3 个随机种子。

\begin{table*}[t]
\centering
\caption{Training and evaluation settings used in the reported experiments.}
\label{tab:training_settings}
\begin{tabularx}{\textwidth}{lcl}
\toprule
Item & Value & Scope \\
\midrule
Parallel environments & 8 & All training runs \\
Rollout length & 128 steps & All training runs \\
PPO epochs & 4 & All training runs \\
Discount factor $\gamma$ & 0.99 & All training runs \\
GAE $\lambda$ & 0.95 & All training runs \\
Clip coefficient & 0.2 & All training runs \\
Entropy coefficient & 0.01 & All training runs \\
Value loss coefficient & 0.5 & All training runs \\
Max gradient norm & 0.5 & All training runs \\
Learning rate & $3\times10^{-4}$ & Default training; staged fine-tuning may use a lower value \\
Hidden dimension & 128 & Actor, critic, and graph encoders \\
Role embedding dimension & 8 & GAT-MAPPO and EA-RG-MAPPO-S \\
Graph minibatch size & 256 graphs & GAT-MAPPO and EA-RG-MAPPO-S \\
MAPPO minibatch size & 512 samples & MAPPO baseline \\
Target policy & mixed & Final evaluation and robustness tests \\
Target speed & 0.75 & Final main evaluation \\
Training communication radius & 8 & Fixed-radius stage \\
Staged fine-tuning radius & $U(4,10)$ & EA-RG-MAPPO-S random-radius stage \\
Final evaluation episodes & 300 per seed & Main table \\
Evaluation seeds & 3 & All reported mean$\pm$std tables \\
Python environment & cac / Python 3.8.20 & Reported runtime \\
PyTorch & 2.4.1+cu124 & Reported runtime \\
GPU & NVIDIA GTX 1650 Ti & Reported runtime \\
\bottomrule
\end{tabularx}
\end{table*}


所有方法均采用集中训练、分散执行设置。MAPPO 使用共享 actor 和集中式 critic；GAT-MAPPO 在此基础上加入节点级图注意力；EA-RG-MAPPO-S 进一步加入角色嵌入、相对边特征和分阶段随机通信半径微调。训练阶段首先在固定通信半径 $R_c=8$ 下学习基础协同策略，然后从固定半径 checkpoint 出发，在 $R_c\sim U(4,10)$ 的随机半径区间内进行短程微调。该设置用于检验策略是否能够适应通信拓扑变化，而不是只适应单一半径。

\subsection{对比方法}

本文比较 MAPPO、GAT-MAPPO 和 EA-RG-MAPPO-S。MAPPO 是无显式图结构的多智能体强化学习基线。GAT-MAPPO 在 MAPPO 基础上加入普通图注意力。EA-RG-MAPPO-S 是本文最终方法，包含边特征增强角色图和分阶段随机半径微调。在消融实验中，还比较 RG-MAPPO 和 EA-RG-MAPPO，用于分析角色图、边特征和 staged 微调的影响。

\subsection{主结果}

\begin{table*}[t]
\centering
\caption{Final 300-episode evaluation under different communication radii. Results are reported as mean$\pm$std over three seeds.}
\label{tab:final_comm_300_results}
\begin{tabular}{lccccc}
\toprule
Method & Radius & Success $\uparrow$ & Collision $\downarrow$ & Timeout $\downarrow$ & Avg. steps $\downarrow$ \\
\midrule
MAPPO & 4 & 0.707$\pm$0.167 & 0.228$\pm$0.099 & 0.072$\pm$0.070 & 84.632$\pm$31.747 \\
MAPPO & 6 & 0.770$\pm$0.145 & 0.218$\pm$0.130 & 0.017$\pm$0.019 & 70.081$\pm$29.733 \\
MAPPO & 8 & 0.810$\pm$0.136 & 0.177$\pm$0.126 & 0.017$\pm$0.015 & 62.143$\pm$25.158 \\
MAPPO & 10 & 0.836$\pm$0.077 & 0.152$\pm$0.067 & 0.014$\pm$0.011 & 58.812$\pm$16.622 \\
\midrule
GAT-MAPPO & 4 & 0.811$\pm$0.026 & 0.136$\pm$0.020 & 0.060$\pm$0.040 & 72.822$\pm$14.039 \\
GAT-MAPPO & 6 & 0.838$\pm$0.051 & 0.126$\pm$0.034 & 0.039$\pm$0.025 & 67.784$\pm$18.906 \\
GAT-MAPPO & 8 & 0.784$\pm$0.054 & 0.179$\pm$0.042 & 0.038$\pm$0.042 & 71.418$\pm$18.202 \\
GAT-MAPPO & 10 & 0.802$\pm$0.037 & 0.156$\pm$0.024 & 0.043$\pm$0.034 & 71.886$\pm$12.111 \\
\midrule
EA-RG-MAPPO-S & 4 & 0.926$\pm$0.004 & 0.054$\pm$0.007 & 0.022$\pm$0.004 & 67.510$\pm$0.631 \\
EA-RG-MAPPO-S & 6 & 0.919$\pm$0.012 & 0.064$\pm$0.006 & 0.018$\pm$0.006 & 64.286$\pm$1.610 \\
EA-RG-MAPPO-S & 8 & 0.890$\pm$0.021 & 0.083$\pm$0.012 & 0.027$\pm$0.014 & 66.790$\pm$4.714 \\
EA-RG-MAPPO-S & 10 & 0.879$\pm$0.017 & 0.086$\pm$0.020 & 0.039$\pm$0.011 & 69.144$\pm$1.739 \\
\bottomrule
\end{tabular}
\end{table*}


最终主图如图~\ref{fig:final_success} 和图~\ref{fig:final_collision} 所示。EA-RG-MAPPO-S 在四个通信半径下取得 0.879--0.926 的成功率，并将碰撞率控制在 0.054--0.086 之间。相比 MAPPO，EA-RG-MAPPO-S 在低通信半径下从数值上明显降低碰撞率。例如 radius=4 时，MAPPO 碰撞率为 $0.228\pm0.099$，而 EA-RG-MAPPO-S 为 $0.054\pm0.007$；在 radius=6 和 radius=8 下，EA-RG-MAPPO-S 的成功率/碰撞率分别为 $0.919/0.064$ 和 $0.890/0.083$。相比 GAT-MAPPO，EA-RG-MAPPO-S 在 radius=8 和 radius=10 下成功率更高、碰撞率更低，并且标准差更小。该结果说明，普通图注意力并不能自动保证跨通信半径稳定性，边特征和随机半径微调共同提供了更稳健的协同表示。

\begin{figure}[t]
    \centering
    \includegraphics[width=0.48\textwidth]{final_300_success_rate.png}
    \caption{300 回合复评下不同通信半径的成功率。}
    \label{fig:final_success}
\end{figure}

\begin{figure}[t]
    \centering
    \includegraphics[width=0.48\textwidth]{final_300_collision_rate.png}
    \caption{300 回合复评下不同通信半径的碰撞率。}
    \label{fig:final_collision}
\end{figure}

\subsection{消融分析}

\begin{table*}[t]
\centering
\caption{Ablation study of role graph, edge-aware attention, and staged random-radius fine-tuning. Results are reported as mean$\pm$std over three seeds.}
\label{tab:ablation_results}
\begin{tabular}{lccccc}
\toprule
Method & Radius & Success $\uparrow$ & Collision $\downarrow$ & Timeout $\downarrow$ & Avg. steps $\downarrow$ \\
\midrule
RG-MAPPO & 4 & 0.907$\pm$0.042 & 0.080$\pm$0.044 & 0.017$\pm$0.006 & 73.200$\pm$7.600 \\
RG-MAPPO & 6 & 0.917$\pm$0.076 & 0.063$\pm$0.059 & 0.023$\pm$0.023 & 72.770$\pm$11.410 \\
RG-MAPPO & 8 & 0.887$\pm$0.146 & 0.097$\pm$0.125 & 0.017$\pm$0.021 & 72.000$\pm$13.450 \\
RG-MAPPO & 10 & 0.907$\pm$0.071 & 0.073$\pm$0.059 & 0.020$\pm$0.017 & 72.040$\pm$10.180 \\
\midrule
EA-RG-MAPPO & 4 & 0.927$\pm$0.021 & 0.037$\pm$0.015 & 0.037$\pm$0.006 & 65.850$\pm$2.780 \\
EA-RG-MAPPO & 6 & 0.887$\pm$0.015 & 0.073$\pm$0.032 & 0.040$\pm$0.036 & 63.410$\pm$2.200 \\
EA-RG-MAPPO & 8 & 0.900$\pm$0.036 & 0.083$\pm$0.012 & 0.020$\pm$0.026 & 63.780$\pm$4.580 \\
EA-RG-MAPPO & 10 & 0.837$\pm$0.032 & 0.113$\pm$0.015 & 0.050$\pm$0.017 & 73.360$\pm$4.120 \\
\midrule
EA-RG-MAPPO-S & 4 & 0.907$\pm$0.012 & 0.067$\pm$0.012 & 0.027$\pm$0.012 & 69.470$\pm$2.300 \\
EA-RG-MAPPO-S & 6 & 0.907$\pm$0.015 & 0.073$\pm$0.021 & 0.020$\pm$0.010 & 62.680$\pm$6.510 \\
EA-RG-MAPPO-S & 8 & 0.883$\pm$0.051 & 0.083$\pm$0.031 & 0.033$\pm$0.032 & 65.210$\pm$7.700 \\
EA-RG-MAPPO-S & 10 & 0.880$\pm$0.020 & 0.090$\pm$0.026 & 0.033$\pm$0.021 & 70.350$\pm$3.540 \\
\bottomrule
\end{tabular}
\end{table*}


100-episode 全消融表保留用于分析模块贡献。RG-MAPPO、EA-RG-MAPPO 和 EA-RG-MAPPO-S 的比较表明，相对边特征能够降低低通信半径下的碰撞率，并缓解部分种子在 radius=8 下的不稳定；分阶段随机通信半径微调能够恢复 radius=10 泛化，使模型在多个通信半径下表现更均衡。由于该表的评估回合数与最终主表不同，本文将其作为模块分析证据，而不与 300 回合主结果混合解释。

\subsection{可视化分析}

per-seed 散点图显示，MAPPO 在不同种子间波动较大，而 EA-RG-MAPPO-S 的成功率和碰撞率更集中。轨迹案例图显示，在相同环境种子下，基线方法可能出现碰撞，而 EA-RG-MAPPO-S 能保持成功追击。注意力热力图显示，通信半径变化会改变图注意力分布，支持有限通信图建模的解释。

\subsection{目标意图分支诊断}

当前实现曾包含目标意图辅助分支，但诊断结果显示 plain accuracy 为 0.587，balanced accuracy 为 0.200（即 balanced accuracy = 0.200），说明该分支主要预测多数类，不能作为高精度意图识别模块。这不能作为强主张。因此，本文不将意图识别作为主贡献。未来若继续强化该方向，需要引入短时历史、目标转弯率等可观测运动特征，并使用 balanced accuracy 作为主要评价指标。
